\subsection{Zahlenmengen}
\label{subsec:fundamentals_number_sets}

Ausgehend von den natürlichen Zahlen werden die Zahlenmengen schrittweise
erweitert. Der Grund ist stets derselbe: Eine Rechenoperation soll ausführbar
sein, ihr Ergebnis liegt aber nicht mehr in der bisherigen Zahlenmenge.

\begin{DefinitionBox}{Abgeschlossenheit}{fundamentals_closure}
  Eine Zahlenmenge $M$ heißt bezüglich einer Rechenoperation $\circ$
  \emph{abgeschlossen} \EnglishTerm{closed}, wenn für alle $a,b\in M$ auch
  \[
    a\circ b\in M
  \]
  gilt, sofern die Rechenoperation definiert ist.
\end{DefinitionBox}

\subsubsection{Natürliche Zahlen}
\label{subsubsec:fundamentals_natural_numbers}

Die \emph{natürlichen Zahlen} dienen dem Zählen:
\begin{align*}
  \N
    &= \Cloze{7cm}{\{1,2,3,4,\ldots\}},\\[0.4em]
  \N_0
    &= \Cloze{7cm}{\N\cup\{0\}=\{0,1,2,3,\ldots\}}.
\end{align*}
Die Addition und Multiplikation natürlicher Zahlen führen wieder zu einer
natürlichen Zahl. Die natürlichen Zahlen sind bezüglich $+$ und $\cdot$
abgeschlossen:
\[
  3+5=\Cloze{1.0cm}{8}\in\N
  \qquad
  3\cdot5=\Cloze{1.2cm}{15}\in\N
  \qquad \text{mit } 3\in\N,\ 5\in\N.
\]

Bezüglich der Subtraktion sind die natürlichen Zahlen dagegen nicht
abgeschlossen:
\[
  3-5=\Cloze{1.2cm}{-2}\notin\N
  \qquad \text{mit } 3\in\N,\ 5\in\N.
\]


\subsubsection{Ganze Zahlen}
\label{subsubsec:fundamentals_integer_numbers}

Damit auch solche Differenzen erfasst werden, ergänzt man die negativen Zahlen
und erhält die \emph{ganzen Zahlen}:
\[
  \Z=\Cloze{9cm}{\{\ldots,-3,-2,-1,0,1,2,3,\ldots\}}.
\]
\begin{center}
  \begingroup
    \def\NumberSetStage{2}
    \CachedTikzFigure{content/resources/chapters/fundamentals/tikz/number_set_inclusions.tex}
  \endgroup
\end{center}

In $\Z$ sind Addition, Subtraktion und Multiplikation abgeschlossen.

\[
  -3+5=\Cloze{1.0cm}{2}\in\Z
  \qquad
  -3-5=\Cloze{1.2cm}{-8}\in\Z
  \qquad
  (-3)\cdot5=\Cloze{1.3cm}{-15}\in\Z
  \qquad \text{mit } -3\in\Z,\ 5\in\Z.
\]

Bezüglich der Division sind die ganzen Zahlen nicht abgeschlossen:
\[
  -3:5=-\frac35=\Cloze{1.8cm}{-0{,}6}\notin\Z
  \qquad \text{mit } -3\in\Z,\ 5\in\Z.
\]


\subsubsection{Rationale Zahlen}
\label{subsubsec:fundamentals_rational_numbers}

Damit auch Quotienten ganzer Zahlen erfasst werden, erweitert man $\Z$ zu den
\emph{rationalen Zahlen}:
\[
  \Q=\left\{\frac pq\,\middle|\,p\in\Z,\ q\in\N\right\}
  \qquad \text{„Menge aller Brüche aus ganzen Zahlen“}.
\]
\begin{center}
  \begingroup
    \def\NumberSetStage{3}
    \CachedTikzFigure{content/resources/chapters/fundamentals/tikz/number_set_inclusions.tex}
  \endgroup
\end{center}
Jede rationale Zahl lässt sich als Bruch ganzer Zahlen darstellen. Endliche
und periodische Dezimalzahlen sind daher rational, zum Beispiel
\[
  -\frac12=\Cloze{1.8cm}{-0{,}5}\in\Q,
  \qquad
  \frac19=\Cloze{2.0cm}{0{,}\overline{1}}\in\Q.
\]
Die rationalen Zahlen sind bezüglich $+$, $-$, $\cdot$ und der Division durch
eine von null verschiedene Zahl abgeschlossen.

\[
  \begin{aligned}
    \frac12+\left(-\frac14\right)
      &=\Cloze{3.0cm}{\frac14}\in\Q,
    &\qquad
    \frac12-\left(-\frac14\right)
      &=\Cloze{3.0cm}{\frac34}\in\Q,\\[0.5em]
    \frac12\cdot\left(-\frac14\right)
      &=\Cloze{3.0cm}{-\frac18}\in\Q,
    &\qquad
    \frac12:\left(-\frac14\right)
      &=\Cloze{3.0cm}{-2}\in\Q
  \end{aligned}
  \qquad
  \text{mit } \frac12\in\Q,\ -\frac14\in\Q.
\]

Bezüglich des Wurzelziehens sind die rationalen Zahlen allerdings nicht
abgeschlossen. So gibt es keine rationale Zahl, deren Quadrat gleich $2$ ist:
\[
  x^2=2
  \quad\Longrightarrow\quad
  x=\pm\Cloze{1.2cm}{\sqrt2}\notin\Q.
\]


\subsubsection{Reelle Zahlen}
\label{subsubsec:fundamentals_real_numbers}

Zahlen, die sich nicht als Quotient ganzer Zahlen darstellen lassen, heißen
\emph{irrationale Zahlen}. Beispiele sind
\[
  \sqrt2=1{,}4142\ldots,
  \qquad
  \pi=3{,}14159\ldots,
  \qquad
  \mathrm e=2{,}71828\ldots.
\]
Rationale und irrationale Zahlen bilden zusammen die \emph{reellen Zahlen}.

\begin{center}
  \begingroup
    \def\NumberSetStage{4}
    \CachedTikzFigure{content/resources/chapters/fundamentals/tikz/number_set_inclusions.tex}
  \endgroup
\end{center}

Die reellen Zahlen $\R$ sind bezüglich Addition, Subtraktion und Multiplikation
sowie bezüglich der Division durch eine von null verschiedene Zahl
abgeschlossen.

\begin{KeyPointBox}{Vollständigkeit der reellen Zahlen}{fundamentals_real_number_completeness}
  Jeder Punkt auf der Zahlengeraden entspricht genau einer reellen Zahl und
  umgekehrt. Anschaulich besitzt die Zahlengerade daher keine Lücken: Die
  reellen Zahlen $\R$ sind \emph{vollständig}.
  \begin{center}
    \CachedTikzFigure{content/resources/chapters/fundamentals/tikz/real_number_line.tex}
  \end{center}
   Damit bilden sie eine zentrale Zahlenmenge in den Naturwissenschaften und der Technik.
\end{KeyPointBox}

Bezüglich des Wurzelziehens sind allerdings auch die reellen Zahlen nicht abgeschlossen:
Die reelle Gleichung
\[
  x^2=-1
\]
besitzt in $\R$ keine Lösung, denn das Quadrat jeder reellen Zahl ist
nichtnegativ.



\subsubsection{Ausblick: Komplexe Zahlen}
\label{subsubsec:fundamentals_complex_numbers}

Um dennoch die Wurzel aus $-1$ ziehen zu können, führt man die imaginäre
Einheit $i$ mit der Eigenschaft
\[
  i^2=-1
\]
ein. Da $i\notin\R$ gilt, ist $i$ keine reelle Zahl. Auf diese Weise erweitert
man die reellen Zahlen $\R$ zu den sogenannten \emph{komplexen Zahlen} $\C$.

\begin{center}
  \begingroup
    \def\NumberSetStage{5}
    \CachedTikzFigure{content/resources/chapters/fundamentals/tikz/number_set_inclusions.tex}
  \endgroup
\end{center}

Eine ausführliche Einführung in die komplexen Zahlen $\C$ folgt in
Kapitel~\ref{chap:complex_numbers}.


